% =====================================================================
% Rubric Deliverable C1 / P3 — Novelty & AI Justification
% MRI-ReportGenerator (EECE503N). Self-contained; compile with tectonic.
% =====================================================================
\documentclass[11pt]{article}
\usepackage[margin=1in]{geometry}
\usepackage{booktabs}
\usepackage{tabularx}
\usepackage{array}
\usepackage{enumitem}
\usepackage{xcolor}
\usepackage{verbatim}
\usepackage{textcomp}
\usepackage[hidelinks]{hyperref}

\newcolumntype{Y}{>{\raggedright\arraybackslash}X}
\setlist{nosep, leftmargin=1.4em}
\sloppy

\title{\textbf{Novelty \& AI Justification}\\[2pt]
\large Rubric Deliverable C1 / P3 --- MRI-ReportGenerator (Cervical Spine MRI Analysis)}
\author{EECE503N Final Project}
\date{2026-06-09}

\begin{document}
\maketitle

\noindent\textbf{Scope.} This document states what is \emph{novel} about MRI-ReportGenerator and
\emph{justifies} the AI design choices behind it. It is the C1/P3 deliverable (novelty /
AI-justification). It deliberately does \textbf{not} re-explain the measurement algorithms or the
validation numbers --- those are documented in the technical-depth material, the full manuscript, the
final validation summaries, and the development journal. This deliverable \emph{references} that body of
evidence and argues the contribution.

\medskip
\noindent\textbf{Medical framing (hard rule).} The system produces measurements and \emph{findings
flagged for physician review}, never diagnoses. The novelty claims below are framed accordingly.

% ---------------------------------------------------------------------
\section{The novelty in one sentence}
The contribution is \textbf{not} a new neural network. It is a \emph{system design}: a cervical-spine
MRI pipeline built from \textbf{healthy-anchored, disease-agnostic geometric detectors}, validated by a
\textbf{threshold-crossing / distribution-separation methodology} purpose-built for a domain where
\textbf{no public dataset pairs cervical MRI with per-case expert measurements}, and surfaced through a
\textbf{fully cited threshold catalog} that only ever flags findings for physician review. Each of those
four italicised choices is a deliberate, defensible departure from the obvious approach, and together
they are what a thin ``wrapper over a segmentation model'' does not have.

% ---------------------------------------------------------------------
\section{Why a non-AI baseline exists --- and why automation is justified}
The non-AI baseline is \textbf{manual radiologist measurement}. It is the right thing to improve on for
two cited reasons (full numbers and citations are summarized in the positioning material):
\begin{itemize}
  \item \textbf{Slow.} Cervical MRI reporting runs $\approx 2.7$\,min median / $3.8$\,min mean of
  PACS interaction (Forsberg 2017, PMID 27714473), with a further self-reported $5$--$10$\,min/case of
  manual geometry (Zhu 2024, PMID 38269650); no formal time-motion study isolates the geometry step ---
  itself a documented gap.
  \item \textbf{Unreliable on exactly what we automate.} Inter-observer agreement collapses for the
  structures this pipeline measures: cord AP ICC $0.66$ axial (Grochmal 2018, PMID 29913296),
  cord-compression metrics ICC $0.35$--$0.56$ (Fehlings 2006, PMID 16816769), cervical disc grading on
  the original Pfirrmann scale $\kappa = 0.265$ (Urbanschitz 2021, PMID 34966859), and Cobb angle
  $\sim 0.88$ for a single trained reader but $\sim 0.55$ across mixed readers (Sevin 2025,
  PMID 41011045).
\end{itemize}
A reproducible, deterministic measurement layer is therefore a genuine improvement, not a solution in
search of a problem. This grounds the \emph{Application} positioning.

% ---------------------------------------------------------------------
\section{Four novelty axes (what is genuinely different)}

\subsection{Disease-agnostic detectors anchored on healthy norms}
Most ``AI for spine MRI'' work trains a classifier to recognise a \emph{named pathology} (stenosis,
herniation, myelopathy) from labelled diseased examples. We do the opposite: we measure
\textbf{geometry and signal against a healthy reference}, so a canal of $8$\,mm reads ``narrow,
flagged for review'' \emph{regardless of cause}. The detector is defined by the normal range, not by a
disease class. Consequences that are themselves the novelty:
\begin{itemize}
  \item It generalises beyond the diseases in any training set --- there is no disease label to overfit.
  \item It needs no diseased training cohort, only a healthy reference distribution, which removes the
  train/test leakage failure mode that plagues small-cohort medical classifiers.
  \item Where a published cervical-MRI norm does not exist, the threshold is \emph{cohort-derived} and
  labelled as such (e.g.\ the Ha/Hp compression z-screen and the disc/VB AP-ratio cut --- both stated as
  cohort-calibrated, not literature values, with the underlying mechanism cited: Machino 2021,
  PMID 34098133, for disc widening with degeneration).
\end{itemize}

\subsection{A validation methodology designed for ``no per-case ground truth''}
The standard evaluation --- per-case sensitivity/specificity against a radiologist gold standard --- is
\textbf{impossible on public data} for cervical MRI, because no dataset pairs the images with per-case
expert measurements (established by repeated dataset searches). Rather than fabricate a weak label, we
designed an evaluation that fits the constraint:
\begin{enumerate}
  \item \textbf{Healthy-anchoring (specificity):} confirm healthy anatomy reads inside the normal range
  on a healthy cohort (Spine-Generic).
  \item \textbf{Threshold-crossing / distribution separation (sensitivity proxy):} confirm symptomatic
  anatomy crosses into the abnormal range on a symptomatic cohort (MMCSD, CSM/CSR), via
  distribution-separation statistics (Mann--Whitney), never per-case accuracy.
\end{enumerate}
Designing and defending a sound evaluation when the textbook one is unavailable is a methods
contribution in its own right (paper \S13; J14--J26).

\subsection{The physical-dimension scanner-immunity insight}
A non-obvious finding shaped the whole evaluation: \textbf{physical-dimension (mm) metrics are
scanner-immune} and validate \emph{across} datasets, while \textbf{intensity/ratio metrics are
acquisition-sensitive} and must be validated \emph{within} one dataset. This is why the millimetre canal
/ SAC / cord measurements (G3) separate cleanly cross-dataset (p$=$0.0001) but disc \emph{signal} does
not, and why disc metrics were validated within-MMCSD only. Recognising this distinction --- and
structuring the validation around it --- is novel methodology, not a parameter choice (J16, J23).

\subsection{Cited, auditable assessement --- never a black box, never a diagnosis}
Every assessed finding maps to a \emph{cited} threshold in the G6 catalog,
with provenance, modality caveats, and an honest label
when a threshold is a documented borrow or a cohort-derived value. The output is a status
(within / review / outside) ``flagged for physician review,'' never a diagnosis. A radiologist can check
each number against its source. This auditability is a deliberate departure from end-to-end
diagnostic models and is the appropriate design for a clinical application.

% ---------------------------------------------------------------------
\section{Honesty about negatives \emph{is} part of the novelty}
A credible system says what does \emph{not} work. We report, and do not hide, three negatives that a
results-chasing project would have buried:
\begin{itemize}
  \item \textbf{G4 cervical alignment (Cobb) is NOT a disease discriminator.} On a balanced multi-site
  cohort (26 healthy vs 41 symptomatic) the effect evaporates: $d=0.28$, AUC $0.57$, $p=0.32$. An earlier
  $n=11$ result ($d=0.76$) was a lordosis-biased small sample. We report G4 as a validated
  \emph{measurement}, not a screen (J24, J26).
  \item \textbf{Disc signal and posterior bulge do not discriminate} within the symptomatic cohort
  (AUC $0.50$ each, level-controlled); only the disc/VB AP-ratio carries a modest signal (AUC $0.62$)
  (J23).
  \item \textbf{The compression-fracture abnormal arm is unvalidated} because no labelled cervical-MRI
  compression-fracture dataset exists (documented limitation; VerSe authors note cervical fractures are
  rare and usually non-osteoporotic, PMC8082364).
\end{itemize}
Reporting these as negatives --- rather than tuning a number until it looks clean --- is the same
discipline that makes the positive results (G3 stenosis, p$=$0.0001) trustworthy.

% ---------------------------------------------------------------------
\section{AI justification: why this is the right AI architecture}
The pretrained-and-compose design (three frozen segmentation engines feeding our measurement and
assessement layer) is the defensible choice for a clinical
application:
\begin{itemize}
  \item \textbf{Auditability.} Outputs trace to a measurement and a cited threshold, not an opaque model.
  \item \textbf{No overfitting by construction.} Frozen segmenters + healthy-anchored / cited thresholds
  mean the system cannot have memorised the evaluation cohort.
  \item \textbf{Graceful degradation.} Per-component error contracts mean one failed measurement (e.g.\ a
  missing SPINEPS mask) degrades a single method, not the whole report.
  \item \textbf{The differentiator is the clinical layer.} The measurement validity, the cited
  assessement, and the report --- the parts that decide whether a radiologist trusts the output ---
  are exactly the parts we built and validated.
\end{itemize}

% ---------------------------------------------------------------------
\section{Relation to similarly-scoped work}
This project shares a high-level title space with other cervical-spine MRI efforts. The differentiation
above is stated on its own terms and does not depend on any one comparator: the combination of
disease-agnostic healthy-anchored detectors, a no-per-case-GT validation methodology, the
scanner-immunity insight, and a fully cited review-only assessement layer is the contribution. A
direct, point-by-point head-to-head against a specific comparator project can be appended here once that
project's scope and abstract are available; it is not required to establish the novelty, which stands on
the four axes above.
% NOTE: if a specific comparator (e.g. a duplicate-title team) shares its abstract/scope, add a
% head-to-head differentiation table in this section. The novelty argument above does NOT depend on it.

% ---------------------------------------------------------------------
\section{Summary for the grader}
\begin{itemize}
  \item The novelty is a \textbf{system}: healthy-anchored disease-agnostic detectors, a validation
  methodology built for the no-per-case-GT constraint, the physical-dimension scanner-immunity insight,
  and a cited review-only assessement layer.
  \item The AI choice (compose frozen, validated segmenters; build and validate the clinical layer) is
  \textbf{justified}, not a shortcut --- it gives auditability, no overfitting, and graceful degradation.
  \item Credibility comes from \textbf{reported negatives} (G4 not a discriminator; disc signal/bulge
  dead; compression-fracture arm unvalidated) as much as from the strong G3 result.
  \item Full technical depth, methods, numbers, and development narrative are documented in the
  companion technical materials and the full manuscript.
\end{itemize}

\end{document}
